% =====================================================================
%  Danh mục từ viết tắt (Glossary / List of Abbreviations)
%  Med_Pharm_Bio_Nexus - AI Conquer 2026
%
%  CÁCH DÙNG:
%  1. Bảo đảm trong main.tex (phần preamble) có nạp các gói sau:
%         \usepackage{booktabs}
%         \usepackage{longtable}
%         \usepackage{array}
%     (Nếu template AIConquer2026 đã nạp sẵn thì bỏ qua.)
%  2. Chèn vào vị trí mong muốn (thường ngay sau Mục lục hoặc trong Phụ lục):
%         % =====================================================================
%  Danh mục từ viết tắt (Glossary / List of Abbreviations)
%  Med_Pharm_Bio_Nexus - AI Conquer 2026
%
%  CÁCH DÙNG:
%  1. Bảo đảm trong main.tex (phần preamble) có nạp các gói sau:
%         \usepackage{booktabs}
%         \usepackage{longtable}
%         \usepackage{array}
%     (Nếu template AIConquer2026 đã nạp sẵn thì bỏ qua.)
%  2. Chèn vào vị trí mong muốn (thường ngay sau Mục lục hoặc trong Phụ lục):
%         % =====================================================================
%  Danh mục từ viết tắt (Glossary / List of Abbreviations)
%  Med_Pharm_Bio_Nexus - AI Conquer 2026
%
%  CÁCH DÙNG:
%  1. Bảo đảm trong main.tex (phần preamble) có nạp các gói sau:
%         \usepackage{booktabs}
%         \usepackage{longtable}
%         \usepackage{array}
%     (Nếu template AIConquer2026 đã nạp sẵn thì bỏ qua.)
%  2. Chèn vào vị trí mong muốn (thường ngay sau Mục lục hoặc trong Phụ lục):
%         % =====================================================================
%  Danh mục từ viết tắt (Glossary / List of Abbreviations)
%  Med_Pharm_Bio_Nexus - AI Conquer 2026
%
%  CÁCH DÙNG:
%  1. Bảo đảm trong main.tex (phần preamble) có nạp các gói sau:
%         \usepackage{booktabs}
%         \usepackage{longtable}
%         \usepackage{array}
%     (Nếu template AIConquer2026 đã nạp sẵn thì bỏ qua.)
%  2. Chèn vào vị trí mong muốn (thường ngay sau Mục lục hoặc trong Phụ lục):
%         \input{sections/Abbreviations}
%
%  GHI CHÚ NHẤT QUÁN (khuyến nghị sửa trong thân bài):
%    - Thống nhất "k-NN" (bài đang lẫn "kNN" ở Bảng 3).
%    - "CI" mang hai nghĩa: Confidence Interval (thống kê) và
%      Continuous Integration (dựng PDF). Nên viết đầy đủ ở mỗi ngữ cảnh.
% =====================================================================

\section*{Danh mục từ viết tắt}
\addcontentsline{toc}{section}{Danh mục từ viết tắt}

\renewcommand{\arraystretch}{1.3}

\begingroup
\small
\begin{longtable}{@{}p{0.16\textwidth} p{0.40\textwidth} p{0.38\textwidth}@{}}
\toprule
\textbf{Viết tắt} & \textbf{Dạng đầy đủ} & \textbf{Giải nghĩa} \\
\midrule
\endfirsthead

\multicolumn{3}{@{}l}{\small\itshape (tiếp theo trang trước)}\\
\toprule
\textbf{Viết tắt} & \textbf{Dạng đầy đủ} & \textbf{Giải nghĩa} \\
\midrule
\endhead

\midrule
\multicolumn{3}{r@{}}{\small\itshape (tiếp trang sau)}\\
\endfoot

\bottomrule
\endlastfoot

% ---------- Kiến trúc & phương pháp truy xuất ----------
\multicolumn{3}{@{}l}{\textbf{Kiến trúc \& phương pháp truy xuất}}\\
\addlinespace[2pt]
RAG      & Retrieval-Augmented Generation & Sinh văn bản có tăng cường truy xuất tri thức \\
KB       & Knowledge Base & Cơ sở tri thức bệnh học (49 tài liệu) \\
IR       & Information Retrieval & Truy xuất thông tin \\
BM25     & Best Matching 25 (Okapi BM25) & Mô hình truy xuất từ vựng theo khung xác suất \\
IDF      & Inverse Document Frequency & Tần suất tài liệu nghịch đảo (thành phần trọng số BM25) \\
RRF      & Reciprocal Rank Fusion & Hợp nhất thứ hạng theo nghịch đảo hạng ($k=60$) \\
k-NN     & $k$-Nearest Neighbors & Truy xuất $k$ láng giềng gần nhất (dense retrieval) \\
ANN      & Approximate Nearest Neighbor & Tìm láng giềng gần nhất xấp xỉ \\
HNSW     & Hierarchical Navigable Small World & Thuật toán ANN cho tìm kiếm vector xấp xỉ \\
LLM      & Large Language Model & Mô hình ngôn ngữ lớn \\
NLP      & Natural Language Processing & Xử lý ngôn ngữ tự nhiên \\
\addlinespace[4pt]

% ---------- Mô hình & embedding ----------
\multicolumn{3}{@{}l}{\textbf{Mô hình \& embedding}}\\
\addlinespace[2pt]
BGE      & BAAI General Embedding & Họ mô hình embedding (dùng \texttt{bge-small-en-v1.5}, 384 chiều) \\
BAAI     & Beijing Academy of Artificial Intelligence & Đơn vị phát hành họ mô hình BGE \\
BERT     & Bidirectional Encoder Representations from Transformers & Mô hình nền cho cross-encoder re-ranking \\
DPR      & Dense Passage Retrieval & Mô hình truy xuất dense (tham chiếu related work) \\
SBERT    & Sentence-BERT & Mô hình sinh embedding câu (tham chiếu related work) \\
\addlinespace[4pt]

% ---------- Chỉ số đánh giá ----------
\multicolumn{3}{@{}l}{\textbf{Chỉ số đánh giá}}\\
\addlinespace[2pt]
Hit@$k$  & Hit at $k$ & Tỉ lệ truy vấn có bệnh đúng nằm trong top-$k$ \\
MRR      & Mean Reciprocal Rank & Trung bình nghịch đảo thứ hạng của kết quả đúng \\
NDCG@5   & Normalized Discounted Cumulative Gain at 5 & Chất lượng xếp hạng có trọng số (graded relevance) \\
MSR      & Missed-Severe Ratio & Tỉ lệ bỏ sót bệnh nặng (severity $\le 2$) khỏi top-1 \\
CI       & Confidence Interval & Khoảng tin cậy (bootstrap 95\%) -- ngữ cảnh thống kê \\
QA       & Quality Assurance & Kiểm chứng chất lượng \\
\addlinespace[4pt]

% ---------- Dữ liệu & y khoa ----------
\multicolumn{3}{@{}l}{\textbf{Dữ liệu \& thuật ngữ y khoa}}\\
\addlinespace[2pt]
DDXPlus  & Differential Diagnosis Plus & Bộ dữ liệu chẩn đoán phân biệt (49 bệnh, 223 evidence) \\
ICD-10   & International Classification of Diseases, 10th Revision & Phân loại bệnh quốc tế phiên bản 10 \\
CC-BY    & Creative Commons Attribution & Giấy phép của bộ dữ liệu \\
URTI     & Upper Respiratory Tract Infection & Nhiễm trùng đường hô hấp trên \\
GERD     & Gastroesophageal Reflux Disease & Trào ngược dạ dày--thực quản \\
SLE      & Systemic Lupus Erythematosus & Lupus ban đỏ hệ thống \\
NSTEMI   & Non-ST-Elevation Myocardial Infarction & Nhồi máu cơ tim không ST chênh lên \\
STEMI    & ST-Elevation Myocardial Infarction & Nhồi máu cơ tim có ST chênh lên \\
\addlinespace[4pt]

% ---------- Hạ tầng, công cụ & quản trị ----------
\multicolumn{3}{@{}l}{\textbf{Hạ tầng, công cụ \& quản trị dự án}}\\
\addlinespace[2pt]
API      & Application Programming Interface & Giao diện lập trình ứng dụng (FastAPI ở Phase 2) \\
CI       & Continuous Integration & Tích hợp liên tục (GitHub Actions dựng PDF) -- ngữ cảnh DevOps \\
PR       & Pull Request & Yêu cầu gộp mã trên GitHub \\
ADR      & Architecture Decision Record & Bản ghi quyết định kiến trúc (lưu trên Notion) \\
CRediT   & Contributor Roles Taxonomy & Chuẩn phân loại vai trò đóng góp \\
MVP      & Minimum Viable Product & Sản phẩm khả dụng tối thiểu (Phase 1) \\

\end{longtable}
\endgroup

% Ghi chú cuối bảng (tuỳ chọn - có thể xoá nếu không cần)
\vspace{-4pt}
{\footnotesize\itshape
Ghi chú: ``CI'' xuất hiện với hai nghĩa tuỳ ngữ cảnh (Confidence Interval trong phần thống kê;
Continuous Integration trong phần vận hành nhóm). Các định danh mô hình, trường dữ liệu và
lệnh (ví dụ \texttt{keyword\_text}, \texttt{knn\_vector}, \texttt{bge-small-en-v1.5}) được giữ nguyên dạng gốc.
\par}

%
%  GHI CHÚ NHẤT QUÁN (khuyến nghị sửa trong thân bài):
%    - Thống nhất "k-NN" (bài đang lẫn "kNN" ở Bảng 3).
%    - "CI" mang hai nghĩa: Confidence Interval (thống kê) và
%      Continuous Integration (dựng PDF). Nên viết đầy đủ ở mỗi ngữ cảnh.
% =====================================================================

\section*{Danh mục từ viết tắt}
\addcontentsline{toc}{section}{Danh mục từ viết tắt}

\renewcommand{\arraystretch}{1.3}

\begingroup
\small
\begin{longtable}{@{}p{0.16\textwidth} p{0.40\textwidth} p{0.38\textwidth}@{}}
\toprule
\textbf{Viết tắt} & \textbf{Dạng đầy đủ} & \textbf{Giải nghĩa} \\
\midrule
\endfirsthead

\multicolumn{3}{@{}l}{\small\itshape (tiếp theo trang trước)}\\
\toprule
\textbf{Viết tắt} & \textbf{Dạng đầy đủ} & \textbf{Giải nghĩa} \\
\midrule
\endhead

\midrule
\multicolumn{3}{r@{}}{\small\itshape (tiếp trang sau)}\\
\endfoot

\bottomrule
\endlastfoot

% ---------- Kiến trúc & phương pháp truy xuất ----------
\multicolumn{3}{@{}l}{\textbf{Kiến trúc \& phương pháp truy xuất}}\\
\addlinespace[2pt]
RAG      & Retrieval-Augmented Generation & Sinh văn bản có tăng cường truy xuất tri thức \\
KB       & Knowledge Base & Cơ sở tri thức bệnh học (49 tài liệu) \\
IR       & Information Retrieval & Truy xuất thông tin \\
BM25     & Best Matching 25 (Okapi BM25) & Mô hình truy xuất từ vựng theo khung xác suất \\
IDF      & Inverse Document Frequency & Tần suất tài liệu nghịch đảo (thành phần trọng số BM25) \\
RRF      & Reciprocal Rank Fusion & Hợp nhất thứ hạng theo nghịch đảo hạng ($k=60$) \\
k-NN     & $k$-Nearest Neighbors & Truy xuất $k$ láng giềng gần nhất (dense retrieval) \\
ANN      & Approximate Nearest Neighbor & Tìm láng giềng gần nhất xấp xỉ \\
HNSW     & Hierarchical Navigable Small World & Thuật toán ANN cho tìm kiếm vector xấp xỉ \\
LLM      & Large Language Model & Mô hình ngôn ngữ lớn \\
NLP      & Natural Language Processing & Xử lý ngôn ngữ tự nhiên \\
\addlinespace[4pt]

% ---------- Mô hình & embedding ----------
\multicolumn{3}{@{}l}{\textbf{Mô hình \& embedding}}\\
\addlinespace[2pt]
BGE      & BAAI General Embedding & Họ mô hình embedding (dùng \texttt{bge-small-en-v1.5}, 384 chiều) \\
BAAI     & Beijing Academy of Artificial Intelligence & Đơn vị phát hành họ mô hình BGE \\
BERT     & Bidirectional Encoder Representations from Transformers & Mô hình nền cho cross-encoder re-ranking \\
DPR      & Dense Passage Retrieval & Mô hình truy xuất dense (tham chiếu related work) \\
SBERT    & Sentence-BERT & Mô hình sinh embedding câu (tham chiếu related work) \\
\addlinespace[4pt]

% ---------- Chỉ số đánh giá ----------
\multicolumn{3}{@{}l}{\textbf{Chỉ số đánh giá}}\\
\addlinespace[2pt]
Hit@$k$  & Hit at $k$ & Tỉ lệ truy vấn có bệnh đúng nằm trong top-$k$ \\
MRR      & Mean Reciprocal Rank & Trung bình nghịch đảo thứ hạng của kết quả đúng \\
NDCG@5   & Normalized Discounted Cumulative Gain at 5 & Chất lượng xếp hạng có trọng số (graded relevance) \\
MSR      & Missed-Severe Ratio & Tỉ lệ bỏ sót bệnh nặng (severity $\le 2$) khỏi top-1 \\
CI       & Confidence Interval & Khoảng tin cậy (bootstrap 95\%) -- ngữ cảnh thống kê \\
QA       & Quality Assurance & Kiểm chứng chất lượng \\
\addlinespace[4pt]

% ---------- Dữ liệu & y khoa ----------
\multicolumn{3}{@{}l}{\textbf{Dữ liệu \& thuật ngữ y khoa}}\\
\addlinespace[2pt]
DDXPlus  & Differential Diagnosis Plus & Bộ dữ liệu chẩn đoán phân biệt (49 bệnh, 223 evidence) \\
ICD-10   & International Classification of Diseases, 10th Revision & Phân loại bệnh quốc tế phiên bản 10 \\
CC-BY    & Creative Commons Attribution & Giấy phép của bộ dữ liệu \\
URTI     & Upper Respiratory Tract Infection & Nhiễm trùng đường hô hấp trên \\
GERD     & Gastroesophageal Reflux Disease & Trào ngược dạ dày--thực quản \\
SLE      & Systemic Lupus Erythematosus & Lupus ban đỏ hệ thống \\
NSTEMI   & Non-ST-Elevation Myocardial Infarction & Nhồi máu cơ tim không ST chênh lên \\
STEMI    & ST-Elevation Myocardial Infarction & Nhồi máu cơ tim có ST chênh lên \\
\addlinespace[4pt]

% ---------- Hạ tầng, công cụ & quản trị ----------
\multicolumn{3}{@{}l}{\textbf{Hạ tầng, công cụ \& quản trị dự án}}\\
\addlinespace[2pt]
API      & Application Programming Interface & Giao diện lập trình ứng dụng (FastAPI ở Phase 2) \\
CI       & Continuous Integration & Tích hợp liên tục (GitHub Actions dựng PDF) -- ngữ cảnh DevOps \\
PR       & Pull Request & Yêu cầu gộp mã trên GitHub \\
ADR      & Architecture Decision Record & Bản ghi quyết định kiến trúc (lưu trên Notion) \\
CRediT   & Contributor Roles Taxonomy & Chuẩn phân loại vai trò đóng góp \\
MVP      & Minimum Viable Product & Sản phẩm khả dụng tối thiểu (Phase 1) \\

\end{longtable}
\endgroup

% Ghi chú cuối bảng (tuỳ chọn - có thể xoá nếu không cần)
\vspace{-4pt}
{\footnotesize\itshape
Ghi chú: ``CI'' xuất hiện với hai nghĩa tuỳ ngữ cảnh (Confidence Interval trong phần thống kê;
Continuous Integration trong phần vận hành nhóm). Các định danh mô hình, trường dữ liệu và
lệnh (ví dụ \texttt{keyword\_text}, \texttt{knn\_vector}, \texttt{bge-small-en-v1.5}) được giữ nguyên dạng gốc.
\par}

%
%  GHI CHÚ NHẤT QUÁN (khuyến nghị sửa trong thân bài):
%    - Thống nhất "k-NN" (bài đang lẫn "kNN" ở Bảng 3).
%    - "CI" mang hai nghĩa: Confidence Interval (thống kê) và
%      Continuous Integration (dựng PDF). Nên viết đầy đủ ở mỗi ngữ cảnh.
% =====================================================================

\section*{Danh mục từ viết tắt}
\addcontentsline{toc}{section}{Danh mục từ viết tắt}

\renewcommand{\arraystretch}{1.3}

\begingroup
\small
\begin{longtable}{@{}p{0.16\textwidth} p{0.40\textwidth} p{0.38\textwidth}@{}}
\toprule
\textbf{Viết tắt} & \textbf{Dạng đầy đủ} & \textbf{Giải nghĩa} \\
\midrule
\endfirsthead

\multicolumn{3}{@{}l}{\small\itshape (tiếp theo trang trước)}\\
\toprule
\textbf{Viết tắt} & \textbf{Dạng đầy đủ} & \textbf{Giải nghĩa} \\
\midrule
\endhead

\midrule
\multicolumn{3}{r@{}}{\small\itshape (tiếp trang sau)}\\
\endfoot

\bottomrule
\endlastfoot

% ---------- Kiến trúc & phương pháp truy xuất ----------
\multicolumn{3}{@{}l}{\textbf{Kiến trúc \& phương pháp truy xuất}}\\
\addlinespace[2pt]
RAG      & Retrieval-Augmented Generation & Sinh văn bản có tăng cường truy xuất tri thức \\
KB       & Knowledge Base & Cơ sở tri thức bệnh học (49 tài liệu) \\
IR       & Information Retrieval & Truy xuất thông tin \\
BM25     & Best Matching 25 (Okapi BM25) & Mô hình truy xuất từ vựng theo khung xác suất \\
IDF      & Inverse Document Frequency & Tần suất tài liệu nghịch đảo (thành phần trọng số BM25) \\
RRF      & Reciprocal Rank Fusion & Hợp nhất thứ hạng theo nghịch đảo hạng ($k=60$) \\
k-NN     & $k$-Nearest Neighbors & Truy xuất $k$ láng giềng gần nhất (dense retrieval) \\
ANN      & Approximate Nearest Neighbor & Tìm láng giềng gần nhất xấp xỉ \\
HNSW     & Hierarchical Navigable Small World & Thuật toán ANN cho tìm kiếm vector xấp xỉ \\
LLM      & Large Language Model & Mô hình ngôn ngữ lớn \\
NLP      & Natural Language Processing & Xử lý ngôn ngữ tự nhiên \\
\addlinespace[4pt]

% ---------- Mô hình & embedding ----------
\multicolumn{3}{@{}l}{\textbf{Mô hình \& embedding}}\\
\addlinespace[2pt]
BGE      & BAAI General Embedding & Họ mô hình embedding (dùng \texttt{bge-small-en-v1.5}, 384 chiều) \\
BAAI     & Beijing Academy of Artificial Intelligence & Đơn vị phát hành họ mô hình BGE \\
BERT     & Bidirectional Encoder Representations from Transformers & Mô hình nền cho cross-encoder re-ranking \\
DPR      & Dense Passage Retrieval & Mô hình truy xuất dense (tham chiếu related work) \\
SBERT    & Sentence-BERT & Mô hình sinh embedding câu (tham chiếu related work) \\
\addlinespace[4pt]

% ---------- Chỉ số đánh giá ----------
\multicolumn{3}{@{}l}{\textbf{Chỉ số đánh giá}}\\
\addlinespace[2pt]
Hit@$k$  & Hit at $k$ & Tỉ lệ truy vấn có bệnh đúng nằm trong top-$k$ \\
MRR      & Mean Reciprocal Rank & Trung bình nghịch đảo thứ hạng của kết quả đúng \\
NDCG@5   & Normalized Discounted Cumulative Gain at 5 & Chất lượng xếp hạng có trọng số (graded relevance) \\
MSR      & Missed-Severe Ratio & Tỉ lệ bỏ sót bệnh nặng (severity $\le 2$) khỏi top-1 \\
CI       & Confidence Interval & Khoảng tin cậy (bootstrap 95\%) -- ngữ cảnh thống kê \\
QA       & Quality Assurance & Kiểm chứng chất lượng \\
\addlinespace[4pt]

% ---------- Dữ liệu & y khoa ----------
\multicolumn{3}{@{}l}{\textbf{Dữ liệu \& thuật ngữ y khoa}}\\
\addlinespace[2pt]
DDXPlus  & Differential Diagnosis Plus & Bộ dữ liệu chẩn đoán phân biệt (49 bệnh, 223 evidence) \\
ICD-10   & International Classification of Diseases, 10th Revision & Phân loại bệnh quốc tế phiên bản 10 \\
CC-BY    & Creative Commons Attribution & Giấy phép của bộ dữ liệu \\
URTI     & Upper Respiratory Tract Infection & Nhiễm trùng đường hô hấp trên \\
GERD     & Gastroesophageal Reflux Disease & Trào ngược dạ dày--thực quản \\
SLE      & Systemic Lupus Erythematosus & Lupus ban đỏ hệ thống \\
NSTEMI   & Non-ST-Elevation Myocardial Infarction & Nhồi máu cơ tim không ST chênh lên \\
STEMI    & ST-Elevation Myocardial Infarction & Nhồi máu cơ tim có ST chênh lên \\
\addlinespace[4pt]

% ---------- Hạ tầng, công cụ & quản trị ----------
\multicolumn{3}{@{}l}{\textbf{Hạ tầng, công cụ \& quản trị dự án}}\\
\addlinespace[2pt]
API      & Application Programming Interface & Giao diện lập trình ứng dụng (FastAPI ở Phase 2) \\
CI       & Continuous Integration & Tích hợp liên tục (GitHub Actions dựng PDF) -- ngữ cảnh DevOps \\
PR       & Pull Request & Yêu cầu gộp mã trên GitHub \\
ADR      & Architecture Decision Record & Bản ghi quyết định kiến trúc (lưu trên Notion) \\
CRediT   & Contributor Roles Taxonomy & Chuẩn phân loại vai trò đóng góp \\
MVP      & Minimum Viable Product & Sản phẩm khả dụng tối thiểu (Phase 1) \\

\end{longtable}
\endgroup

% Ghi chú cuối bảng (tuỳ chọn - có thể xoá nếu không cần)
\vspace{-4pt}
{\footnotesize\itshape
Ghi chú: ``CI'' xuất hiện với hai nghĩa tuỳ ngữ cảnh (Confidence Interval trong phần thống kê;
Continuous Integration trong phần vận hành nhóm). Các định danh mô hình, trường dữ liệu và
lệnh (ví dụ \texttt{keyword\_text}, \texttt{knn\_vector}, \texttt{bge-small-en-v1.5}) được giữ nguyên dạng gốc.
\par}

%
%  GHI CHÚ NHẤT QUÁN (khuyến nghị sửa trong thân bài):
%    - Thống nhất "k-NN" (bài đang lẫn "kNN" ở Bảng 3).
%    - "CI" mang hai nghĩa: Confidence Interval (thống kê) và
%      Continuous Integration (dựng PDF). Nên viết đầy đủ ở mỗi ngữ cảnh.
% =====================================================================

\section*{Danh mục từ viết tắt}
\addcontentsline{toc}{section}{Danh mục từ viết tắt}

\renewcommand{\arraystretch}{1.3}

\begingroup
\small
\begin{longtable}{@{}p{0.16\textwidth} p{0.40\textwidth} p{0.38\textwidth}@{}}
\toprule
\textbf{Viết tắt} & \textbf{Dạng đầy đủ} & \textbf{Giải nghĩa} \\
\midrule
\endfirsthead

\multicolumn{3}{@{}l}{\small\itshape (tiếp theo trang trước)}\\
\toprule
\textbf{Viết tắt} & \textbf{Dạng đầy đủ} & \textbf{Giải nghĩa} \\
\midrule
\endhead

\midrule
\multicolumn{3}{r@{}}{\small\itshape (tiếp trang sau)}\\
\endfoot

\bottomrule
\endlastfoot

% ---------- Kiến trúc & phương pháp truy xuất ----------
\multicolumn{3}{@{}l}{\textbf{Kiến trúc \& phương pháp truy xuất}}\\
\addlinespace[2pt]
RAG      & Retrieval-Augmented Generation & Sinh văn bản có tăng cường truy xuất tri thức \\
KB       & Knowledge Base & Cơ sở tri thức bệnh học (49 tài liệu) \\
IR       & Information Retrieval & Truy xuất thông tin \\
BM25     & Best Matching 25 (Okapi BM25) & Mô hình truy xuất từ vựng theo khung xác suất \\
IDF      & Inverse Document Frequency & Tần suất tài liệu nghịch đảo (thành phần trọng số BM25) \\
RRF      & Reciprocal Rank Fusion & Hợp nhất thứ hạng theo nghịch đảo hạng ($k=60$) \\
k-NN     & $k$-Nearest Neighbors & Truy xuất $k$ láng giềng gần nhất (dense retrieval) \\
ANN      & Approximate Nearest Neighbor & Tìm láng giềng gần nhất xấp xỉ \\
HNSW     & Hierarchical Navigable Small World & Thuật toán ANN cho tìm kiếm vector xấp xỉ \\
LLM      & Large Language Model & Mô hình ngôn ngữ lớn \\
NLP      & Natural Language Processing & Xử lý ngôn ngữ tự nhiên \\
\addlinespace[4pt]

% ---------- Mô hình & embedding ----------
\multicolumn{3}{@{}l}{\textbf{Mô hình \& embedding}}\\
\addlinespace[2pt]
BGE      & BAAI General Embedding & Họ mô hình embedding (dùng \texttt{bge-small-en-v1.5}, 384 chiều) \\
BAAI     & Beijing Academy of Artificial Intelligence & Đơn vị phát hành họ mô hình BGE \\
BERT     & Bidirectional Encoder Representations from Transformers & Mô hình nền cho cross-encoder re-ranking \\
DPR      & Dense Passage Retrieval & Mô hình truy xuất dense (tham chiếu related work) \\
SBERT    & Sentence-BERT & Mô hình sinh embedding câu (tham chiếu related work) \\
\addlinespace[4pt]

% ---------- Chỉ số đánh giá ----------
\multicolumn{3}{@{}l}{\textbf{Chỉ số đánh giá}}\\
\addlinespace[2pt]
Hit@$k$  & Hit at $k$ & Tỉ lệ truy vấn có bệnh đúng nằm trong top-$k$ \\
MRR      & Mean Reciprocal Rank & Trung bình nghịch đảo thứ hạng của kết quả đúng \\
NDCG@5   & Normalized Discounted Cumulative Gain at 5 & Chất lượng xếp hạng có trọng số (graded relevance) \\
MSR      & Missed-Severe Ratio & Tỉ lệ bỏ sót bệnh nặng (severity $\le 2$) khỏi top-1 \\
CI       & Confidence Interval & Khoảng tin cậy (bootstrap 95\%) -- ngữ cảnh thống kê \\
QA       & Quality Assurance & Kiểm chứng chất lượng \\
\addlinespace[4pt]

% ---------- Dữ liệu & y khoa ----------
\multicolumn{3}{@{}l}{\textbf{Dữ liệu \& thuật ngữ y khoa}}\\
\addlinespace[2pt]
DDXPlus  & Differential Diagnosis Plus & Bộ dữ liệu chẩn đoán phân biệt (49 bệnh, 223 evidence) \\
ICD-10   & International Classification of Diseases, 10th Revision & Phân loại bệnh quốc tế phiên bản 10 \\
CC-BY    & Creative Commons Attribution & Giấy phép của bộ dữ liệu \\
URTI     & Upper Respiratory Tract Infection & Nhiễm trùng đường hô hấp trên \\
GERD     & Gastroesophageal Reflux Disease & Trào ngược dạ dày--thực quản \\
SLE      & Systemic Lupus Erythematosus & Lupus ban đỏ hệ thống \\
NSTEMI   & Non-ST-Elevation Myocardial Infarction & Nhồi máu cơ tim không ST chênh lên \\
STEMI    & ST-Elevation Myocardial Infarction & Nhồi máu cơ tim có ST chênh lên \\
\addlinespace[4pt]

% ---------- Hạ tầng, công cụ & quản trị ----------
\multicolumn{3}{@{}l}{\textbf{Hạ tầng, công cụ \& quản trị dự án}}\\
\addlinespace[2pt]
API      & Application Programming Interface & Giao diện lập trình ứng dụng (FastAPI ở Phase 2) \\
CI       & Continuous Integration & Tích hợp liên tục (GitHub Actions dựng PDF) -- ngữ cảnh DevOps \\
PR       & Pull Request & Yêu cầu gộp mã trên GitHub \\
ADR      & Architecture Decision Record & Bản ghi quyết định kiến trúc (lưu trên Notion) \\
CRediT   & Contributor Roles Taxonomy & Chuẩn phân loại vai trò đóng góp \\
MVP      & Minimum Viable Product & Sản phẩm khả dụng tối thiểu (Phase 1) \\

\end{longtable}
\endgroup

% Ghi chú cuối bảng (tuỳ chọn - có thể xoá nếu không cần)
\vspace{-4pt}
{\footnotesize\itshape
Ghi chú: ``CI'' xuất hiện với hai nghĩa tuỳ ngữ cảnh (Confidence Interval trong phần thống kê;
Continuous Integration trong phần vận hành nhóm). Các định danh mô hình, trường dữ liệu và
lệnh (ví dụ \texttt{keyword\_text}, \texttt{knn\_vector}, \texttt{bge-small-en-v1.5}) được giữ nguyên dạng gốc.
\par}
